\section{分支与局部连通性}\label{section:ch03-components-locally-connected}

对于任意给定的空间 $X$，有一种自然的方法将它分成一些连通的（或道路连通的）块。现在我们处理这个问题。

\begin{definition}\label{definition:ch03-component}
在给定的空间 $X$ 中定义一个等价关系如下：若 $X$ 中存在包含 $x$ 与 $y$ 的连通子空间，则规定 $x\sim y$。每一等价类称为 $X$ 的一个\textbf{分支}（component）（或连通分支）。
\end{definition}

这个关系的对称性与自反性是显然的。传递性的推导如下：设 $A$ 是包含 $x$ 和 $y$ 的一个连通子空间，并且 $B$ 是包含 $y$ 和 $z$ 的一个连通子空间，因为 $A$ 和 $B$ 有公共点 $y$，所以 $A\cup B$ 是包含 $x$ 和 $z$ 的连通子空间。

对于 $X$ 的分支有以下特点：

\begin{theorem}\label{theorem:ch03-components-partition}
$X$ 的所有分支是 $X$ 中这样一些两两无交的连通子空间，它们的并等于 $X$，并且 $X$ 中的每一个非空的连通子空间仅与一个分支相交。
\end{theorem}

\begin{proof}
作为等价类的全体，$X$ 的所有分支是两两无交的，并且其并等于 $X$。$X$ 的每一个连通子空间 $A$ 仅与一个分支相交。事实上，若 $A$ 与 $X$ 的分支 $C_{1}$ 和 $C_{2}$ 都相交，设分别有交点 $x_{1}$ 和 $x_{2}$，那么根据定义有 $x_{1}\sim x_{2}$，而这是不可能的，除非 $C_{1}=C_{2}$。

为了证明分支 $C$ 是连通的，在 $C$ 中取一点 $x_{0}$。对于 $C$ 的每一个点 $x$ 我们有 $x_{0}\sim x$，从而存在一个包含 $x_{0}$ 和 $x$ 的连通子空间 $A_{x}$。根据刚刚证明的结果可见，$A_{x}\subset C$。因此，
\begin{align*}
C=\bigcup_{x\in C}A_{x}.
\end{align*}
因为这些子空间 $A_{x}$ 都是连通的并且有一个公共点 $x_{0}$，所以它们的并是连通的。
\end{proof}

\begin{definition}\label{definition:ch03-path-component}
在空间 $X$ 上规定另外一种等价关系如下：如果在 $X$ 中存在一个从 $x$ 到 $y$ 的道路，我们规定 $x\sim y$。每一等价类称为 $X$ 的一个\textbf{道路连通分支}（path component）。
\end{definition}

以下证明这是一个等价关系。首先注意，如果存在一个以闭区间 $[a,b]$ 为定义域的从 $x$ 到 $y$ 的道路 $f\colon[a,b]\to X$，则也存在一个以闭区间 $[c,d]$ 为定义域的从 $x$ 到 $y$ 的道路 $g$（因为 $\mathbb{R}$ 中任何两个闭区间 $[a,b]$ 和 $[c,d]$ 都是同胚的）。对于 $X$ 中每一个 $x$，自反性 $x\sim x$ 可由常值道路 $f\colon[a,b]\to X$ 直接得到，其中，对于所有的 $t$ 有 $f(t)=x$。其次，对称性可由以下事实推出：若 $f\colon[0,1]\to X$ 是从 $x$ 到 $y$ 的一条道路，那么由 $g(t)=f(1-t)$ 定义的``逆道路''$g\colon[0,1]\to X$ 就是从 $y$ 到 $x$ 的一条道路。最后，传递性的证明如下：设从 $x$ 到 $y$ 的一条道路为 $f\colon[0,1]\to X$，从 $y$ 到 $z$ 的一条道路为 $g\colon[1,2]\to X$，我们``黏结 $f$ 和 $g$''而得到一个从 $x$ 到 $z$ 的道路 $h\colon[0,2]\to X$，由黏结引理可知，道路 $h$ 是连续的。

与前面一个定理的证明相仿，我们有以下定理：

\begin{theorem}\label{theorem:ch03-path-components-partition}
$X$ 的所有道路分支是 $X$ 中这样一些两两无交的道路连通子空间，它们的并等于 $X$，并且 $X$ 中每一个非空道路连通子空间仅与一个道路分支相交。
\end{theorem}

注意，由于 $X$ 的连通子空间的闭包是连通的，所以空间 $X$ 的每一个分支是 $X$ 中的闭集。如果 $X$ 只有有限多个分支，那么每一个分支也是 $X$ 中的开集，这是由于每一个分支的补是有限个闭集的并。但是，一般情况下，$X$ 的分支未必是 $X$ 的开集。

$X$ 的道路连通分支则不具备上述这些好的性质，它们未必是 $X$ 中的开集，也未必是 $X$ 中的闭集。考虑以下例子。

\begin{example}\label{example:ch03-q-components}
如果 $\mathbb{Q}$ 是所有有理数组成的 $\mathbb{R}$ 的子空间，则 $\mathbb{Q}$ 的每一分支为单点集。$\mathbb{Q}$ 的每一分支都不是 $\mathbb{Q}$ 的开集。
\end{example}

\begin{example}\label{example:ch03-sine-components}
前一节中给出的``拓扑学家的正弦曲线''$\overline{S}$ 是具有一个分支（因为它是连通的）和两个道路分支的空间。其中，一个道路分支为 $S$，另一个道路分支为垂直区间 $V=0\times[-1,1]$。注意到 $S$ 为 $\overline{S}$ 的非闭的开集，而 $V$ 为 $\overline{S}$ 的非开的闭集。

如果从 $\overline{S}$ 中去掉 $V$ 中所有第二坐标为有理数的点，则所得到的空间仅有一个分支但有不可数个道路分支。
\end{example}

对一个空间来说，连通性是一个有用的性质。但是，在某些场合下，空间局部地满足连通性条件则更为重要。粗略地说，局部连通性就是每一个点处都有一个``任意小''的连通邻域。精确定义如下：

\begin{definition}\label{definition:ch03-locally-connected}
空间 $X$ 称为在 $x$ 处\textbf{局部连通的}（locally connected at $x$），如果对于 $x$ 的每一个邻域 $U$，存在 $x$ 的一个连通邻域 $V$ 包含于 $U$。若 $X$ 在它的每一个点处都是局部连通的，则简称 $X$ 是局部连通的（locally connected）。类似地，空间 $X$ 称为在 $x$ 处\textbf{局部道路连通的}（locally path connected at $x$），如果对于 $x$ 的每一个邻域 $U$，都存在 $x$ 的一个道路连通邻域 $V$ 包含于 $U$。若 $X$ 在它的每一点处都是局部道路连通的，则称 $X$ 是局部道路连通的（locally path connected）。
\end{definition}

\begin{example}\label{example:ch03-locally-connected-examples}
实直线中的每一个区间和每一条射线都是连通且局部连通的。$\mathbb{R}$ 的子空间 $[-1,0)\cup(0,1]$ 是不连通的，但它是局部连通的。拓扑学家的正弦曲线是连通的但不是局部连通的。有理数集 $\mathbb{Q}$ 既不连通也不局部连通。
\end{example}

\begin{theorem}\label{theorem:ch03-locally-connected-open-components}
空间 $X$ 是局部连通的当且仅当 $X$ 中的任何一个开集 $U$ 的每一个分支在 $X$ 中都是开的。
\end{theorem}

\begin{proof}
假设 $X$ 是局部连通的，$U$ 是 $X$ 中的一个开集，$C$ 是 $U$ 的一个分支。如果 $x$ 是 $C$ 的点，我们可以选取 $x$ 的一个连通邻域 $V$ 使得 $V\subset U$。因为 $V$ 是连通的，它就必然包含在 $U$ 的分支 $C$ 中。因此，$C$ 在 $X$ 中是开的。

反之，若 $X$ 中开集的分支在 $X$ 中是开的。给定 $X$ 的一个点 $x$ 以及 $x$ 的一个邻域 $U$，令 $C$ 是 $U$ 中包含 $x$ 的分支。因为 $C$ 是连通的，并且根据假设它在 $X$ 中是开的，所以 $X$ 在 $x$ 处局部连通。
\end{proof}

可以类似地证明以下定理：

\begin{theorem}\label{theorem:ch03-locally-path-connected-open-components}
空间 $X$ 是局部道路连通的当且仅当 $X$ 中的任何一个开集 $U$ 的每一个道路分支在 $X$ 中都是开的。
\end{theorem}

以下定理给出了道路分支与分支之间的关系：

\begin{theorem}\label{theorem:ch03-path-component-in-component}
若 $X$ 是一个拓扑空间，则 $X$ 的每一个道路分支必定包含在 $X$ 的一个分支之中。若 $X$ 是局部道路连通的，则 $X$ 的分支和道路分支相同。
\end{theorem}

\begin{proof}
设 $C$ 是 $X$ 的一个分支，$x$ 是 $C$ 的一个点，$P$ 是 $X$ 的包含 $x$ 的那个道路分支。因为 $P$ 是连通的，所以 $P\subset C$。如果 $X$ 是局部道路连通的，我们要证明 $P=C$。假设 $P\subsetneq C$。令 $Q$ 表示 $X$ 中所有不同于 $P$ 并且都与 $C$ 相交的那些道路分支的并，其中每一个必定都要包含于 $C$，于是
\begin{align*}
C=P\cup Q.
\end{align*}
因为 $X$ 是局部道路连通的，所以 $X$ 的每一个道路分支在 $X$ 中都是开的。因此，$P$（它是一个道路分支）和 $Q$（它是道路分支的并）在 $X$ 中都是开的，它们构成 $C$ 的一个分割。这与 $C$ 的连通性矛盾。
\end{proof}

\subsection*{习题}

\begin{enumerate}
\item $\mathbb{R}_{\ell}$ 的分支和道路分支是什么？有哪些连续映射 $f\colon\mathbb{R}\to\mathbb{R}_{\ell}$？
\item (a) $\mathbb{R}^{\omega}$（关于积拓扑）的分支和道路分支是什么？
\item (b) 考虑具有一致拓扑的 $\mathbb{R}^{\omega}$，证明：$x$ 和 $y$ 属于 $\mathbb{R}^{\omega}$ 的同一个分支当且仅当序列
\begin{align*}
x-y=(x_{1}-y_{1},x_{2}-y_{2},\cdots)
\end{align*}
是有界的。[提示：只要考虑 $y=\boldsymbol{0}$ 的情形。]
\item (c) 赋予 $\mathbb{R}^{\omega}$ 箱拓扑。证明：$x$ 和 $y$ 属于 $\mathbb{R}^{\omega}$ 的同一分支当且仅当序列 $x-y$``终于零''。[提示：若 $x-y$ 不终于零，证明存在 $\mathbb{R}^{\omega}$ 的自同胚 $h$，使得 $h(x)$ 有界而 $h(y)$ 无界。]
\item 证明有序矩形是局部连通但不是局部道路连通的。这个空间的道路分支是什么？
\item 设 $X$ 是局部道路连通的。证明 $X$ 的每一个连通开集是道路连通的。
\item 设 $X$ 表示 $\mathbb{R}^{2}$ 中的区间 $[0,1]\times 0$ 上的所有有理点。$T$ 表示连接点 $p=0\times 1$ 与 $X$ 中点的所有线段的并。
\begin{enumerate}
\item 证明 $T$ 是道路连通的，并且仅在点 $p$ 处是局部连通的。
\item 求 $\mathbb{R}^{2}$ 的一个子空间，使得它是道路连通的，并且在它的任何点处都不是局部连通的。
\end{enumerate}
\item 若对于 $x\in X$ 的每一个邻域 $U$，存在 $X$ 的一个连通子空间包含于 $U$，并且这个连通子空间包含着 $x$ 的某一个邻域，则称空间 $X$ 在 $x$ 处\textbf{弱局部连通}（weakly locally connected at $x$）。证明，若 $X$ 在它的每一点处都是弱局部连通的，则 $X$ 是局部连通的。[提示：证明开集的分支是开的。]
\item 考虑画在\cref{figure:ch03-infinite-broom}中的那个``无穷扫帚''$X$。证明：$X$ 在点 $p$ 不是局部连通的，但在点 $p$ 是弱局部连通的。[提示：$p$ 的任何连通邻域必定包含着所有点 $a_{i}$。]
\item 设 $p\colon X\to Y$ 是一个商映射。证明：若 $X$ 局部连通，则 $Y$ 局部连通。[提示：若 $C$ 为 $Y$ 的开集 $U$ 的一个分支，证明 $p^{-1}(C)$ 为 $p^{-1}(U)$ 的所有分支的并。]
\item 设 $G$ 是一个拓扑群，$C$ 是 $G$ 的含有单位元 $e$ 的那个分支。证明 $C$ 是 $G$ 的一个正规子群。[提示：若 $x\in G$，则 $xC$ 是 $G$ 含有 $x$ 的那个分支。]
\item 设 $X$ 是一个空间。若不存在 $X$ 的由无交开集 $A$ 和 $B$ 所组成的分割 $X=A\cup B$ 使得 $x\in A$ 和 $y\in B$，我们就规定 $x\sim y$。
\begin{enumerate}
\item 证明这是一个等价关系。其等价类我们称为 $X$ 的\textbf{拟分支}（quasicomponents）。
\item 证明 $X$ 的每一个分支包含在 $X$ 的一个拟分支之中。若 $X$ 是局部连通的，则 $X$ 的分支与拟分支相同。
\item 令 $K$ 表示集合 $\{1/n\mid n\in\mathbb{Z}_{+}\}$，$-K$ 表示集合 $\{-1/n\mid n\in\mathbb{Z}_{+}\}$。试确定 $\mathbb{R}^{2}$ 的下列子空间的分支、道路分支以及拟分支：
\begin{align*}
A=(K\times[0,1])\cup\{0\times 0\}\cup\{0\times 1\}.
\end{align*}
\begin{align*}
B=A\cup([0,1]\times\{0\}).
\end{align*}
\begin{align*}
C=(K\times[0,1])\cup(-K\times[-1,0])\cup([0,1]\times-K)\cup([-1,0]\times K).
\end{align*}
\end{enumerate}
\end{enumerate}

\par
\begin{figure}[htbp]
\centering
\includegraphics[width=0.6\textwidth]{ch03-fig25-1.png}
\caption{无穷扫帚}\label{figure:ch03-infinite-broom}
\end{figure}